%==============================================================================
%  CAPÍTULO — LA NORMATIVA ADMINISTRATIVA DE LA SUPERIR
%==============================================================================
\chapter{La normativa administrativa de la SUPERIR como fuente del derecho concursal}
\label{cap:normativa-superir}

\fichaCapitulo{La ley no basta para tramitar.}{Potestad normativa de la Superintendencia,
jerarquía de sus instrumentos y catálogo sistematizado de las normas de carácter general,
instructivos, oficios de instrucción general, oficios circulares y resoluciones exentas
vigentes.}

%==============================================================================
\bloqueTeorico
%==============================================================================

\subsection{Una fuente que la doctrina suele ignorar}

Quien estudia el derecho concursal chileno únicamente en la ley conoce la mitad del
sistema. La operación cotidiana de los procedimientos ---qué formulario se usa, cómo se
acredita una deuda, cómo se cita a una junta telemática, cómo se emite un finiquito, qué
se informa a la Superintendencia y con qué periodicidad--- está regulada por instrumentos
administrativos que la \LIR{} \citeyearpar{ley-20720} autoriza y que rara vez figuran en
los manuales.

Esta omisión doctrinaria tiene un costo forense directo: un escrito impecable en lo
jurídico puede ser objetado por no ajustarse al modelo aprobado por una norma de carácter
general, y un liquidador diligente puede ser sancionado por no informar un hecho que un
instructivo califica como relevante.

\subsection{Fundamento de la potestad normativa}

La \superir{} no legisla. Su potestad es de ejecución: la ley le encomienda determinar la
forma, el contenido y la oportunidad de actuaciones cuyo contenido sustantivo ella misma
define. La distinción es relevante para fijar el límite del control: una norma
administrativa que exceda ese marco es susceptible de reproche, sea por la vía de la
reclamación administrativa, sea ante la Contraloría General de la República.

\begin{alerta}[El límite de la potestad]
La normativa administrativa puede detallar cómo se cumple un deber legal; no puede crear
deberes nuevos ni restringir derechos que la ley concede. La frontera es difusa y ha sido
objeto de discusión, particularmente respecto de las exigencias documentales impuestas a
los deudores.
\end{alerta}

%==============================================================================
\bloqueNormativo
%==============================================================================

\subsection{Jerarquía y tipología de los instrumentos}

\begin{description}[leftmargin=0pt,style=nextline,font=\sffamily\bfseries]
  \item[Normas de Carácter General (NCG)] Instrumento de mayor alcance. Regulan de modo
  permanente y general una materia entregada a la Superintendencia: modelos de
  presentación, requisitos de las nóminas, forma de las cuentas, mecanismos de
  nominación. Se aprueban por resolución exenta.

  \item[Instructivos] Dirigidos a los sujetos fiscalizados ---veedores, liquidadores,
  síndicos, martilleros---, precisan cómo deben ejecutarse determinadas actuaciones.
  También se aprueban por resolución exenta.

  \item[Oficios de instrucción general] Resuelven cuestiones interpretativas concretas
  con vocación de aplicación general. Constituyen, en los hechos, la doctrina
  administrativa del órgano.

  \item[Oficios circulares] Imparten pautas y recomendaciones de buenas prácticas, con
  menor densidad normativa.

  \item[Resoluciones exentas] Actos administrativos que aprueban, modifican, complementan
  o derogan los instrumentos anteriores, además de regular materias propias de la gestión
  del órgano.
\end{description}

\subsection{Catálogo de Normas de Carácter General vigentes}

\begin{table}[htbp]
\centering
\small
\caption{Normas de Carácter General vigentes}
\label{tab:ncg-vigentes}
\begin{tabularx}{\textwidth}{@{}p{0.09\textwidth} p{0.13\textwidth} X@{}}
\toprule
\textbf{NCG} & \textbf{Fecha} & \textbf{Materia} \\
\midrule
28 & 19.11.2025 & Procedimiento concursal de renegociación de la persona deudora \\
27 & 04.09.2024 & Forma y contenidos de las cuentas provisorias de los liquidadores \\
26 & 07.05.2024 & Requisitos e indicadores para veedores y liquidadores concursales \\
25 & 13.10.2023 & Uso de plataformas electrónicas para la enajenación de bienes muebles \\
24 & 11.08.2023 & Modelo de propuesta de acuerdo en la reorganización simplificada \\
23 & 11.08.2023 & Declaración jurada para la reorganización simplificada de MIPE \\
22 & 11.08.2023 & Declaración jurada para la liquidación voluntaria simplificada \\
20 & 11.08.2023 & Informe sobre cumplimiento del acuerdo de reorganización \\
19 & 11.08.2023 & Certificado de estado de deudas y declaración jurada en reorganización \\
18 & 11.08.2023 & Objeción de la cuenta final de administración \\
17 & 11.08.2023 & Nominación mediante sorteo de veedores o liquidadores \\
16 & 11.08.2023 & Exámenes de conocimientos para veedores, liquidadores y martilleros \\
15 & 11.08.2023 & Realización telemática de las juntas de acreedores \\
14 & 11.08.2023 & Publicaciones en el Boletín Concursal y notificaciones electrónicas \\
10 & 31.12.2019 & Garantía de fiel desempeño \\
7  & 08.10.2014 & Cuentas provisorias y cuenta final de administración \\
6  & 08.10.2014 & Forma de otorgar garantía a los acreedores de primera clase \\
4  & 05.09.2014 & Modelo de solicitud de reorganización \\
\bottomrule
\end{tabularx}
\end{table}

\begin{foco}[Dos lecturas del catálogo]
Primera: la concentración de normas con fecha 11 de agosto de 2023 no es casual. Ese día
entró en vigencia la \Lreforma{} \citeyearpar{ley-21563}, tres meses después de su
publicación, y la Superintendencia dictó simultáneamente el cuerpo normativo que la
operativiza.

Segunda: la ausencia de la NCG \Nro 21 en el catálogo vigente. La norma que regulaba las
pautas de ingreso de la solicitud de renegociación fue sustituida por la \ncg{28}
\citeyearpar{ncg-28}. Toda referencia doctrinaria a la NCG \Nro 21 debe considerarse
superada.
\end{foco}

\subsection{Instructivos vigentes de la Ley \Nro 20.720}

\begin{table}[htbp]
\centering
\small
\caption{Instructivos SUPERIR y resoluciones que los aprueban}
\label{tab:instructivos}
\begin{tabularx}{\textwidth}{@{}p{0.07\textwidth} p{0.13\textwidth} X p{0.17\textwidth}@{}}
\toprule
\textbf{N.\textsuperscript{o}} & \textbf{Fecha} & \textbf{Materia} & \textbf{Aprobado por} \\
\midrule
6 & 04.11.2024 & Finiquitos laborales electrónicos y firma electrónica en la liquidación
& Res. Ex. 16.245 \\
5 & 03.10.2024 & Incautación, conservación y entrega de bienes en liquidación ordinaria y
simplificada & Res. Ex. 14.477 \\
4 & 05.09.2024 & Hechos relevantes que informan los sujetos fiscalizados & Res. Ex. 13.139 \\
3 & 26.08.2024 & Instrucciones a los veedores en la reorganización & Res. Ex. 12.473 \\
2 & 13.08.2024 & Síndicos en el sobreseimiento definitivo y temporal (Libro IV CCom)
& Res. Ex. 11.679 \\
1 & 04.04.2024 & Liquidadores que cesan anticipadamente, suplentes y titulares (art. 38)
& Res. Ex. 4.389 \\
1 & 19.08.2025 & Fianzas, preferencias y contragarantías con sociedades anónimas de
garantía recíproca & Res. Ex. 15.573 \\
1 & 21.01.2026 & Síndicos que cesan anticipadamente y síndicos suplentes & Res. Ex. 999 \\
\bottomrule
\end{tabularx}
\end{table}

\begin{alerta}[Numeración reiterada]
La Superintendencia reinicia la numeración de instructivos por año, de modo que coexisten
varios «Instructivo N.\textsuperscript{o} 1» de contenido distinto. Al citarlos debe
indicarse siempre la resolución aprobatoria y su fecha, y no sólo el número.
\end{alerta}

\subsection{Modificaciones posteriores relevantes}

El Instructivo \Nro 5, sobre incautación, fue objeto de una derogación parcial: la
\resex{16.492}{2025} dejó sin efecto el numeral 3 de su artículo 12
\parencite{resex-16492}. Quien invoque ese instructivo debe verificar la vigencia del
precepto concreto que cita.

\subsection{Oficios de instrucción general}

Constituyen la doctrina administrativa del órgano. Los de mayor incidencia práctica,
agrupados por materia:

\begin{description}[leftmargin=0pt,style=nextline,font=\sffamily\bfseries]
  \item[Materia laboral y previsional] Imputación del aporte al seguro de cesantía en las
  indemnizaciones \parencite{oficio-5080}; pagos y empozamientos ante AFP y AFC
  \parencite{oficio-334}; entrada en vigencia de modificaciones al DL \Nro 3.500 y a la
  Ley \Nro 19.728 \parencite{oficio-15795}; procedencia del pago o reserva de créditos
  laborales \parencite{oficio-1765, oficio-6499}.

  \item[Realización de activos] Reemplazo de bienes muebles en la liquidación
  simplificada \parencite{oficio-9378}; información para la incautación de acreencias
  bancarias \parencite{oficio-10657}.

  \item[Administración de fondos] Pagos pendientes de cobro y empozamientos
  \parencite{oficio-5003}; apertura de cuentas corrientes e inversión de fondos
  \parencite{oficio-11628}.

  \item[Materia tributaria] Obligaciones tributarias de los deudores en liquidación
  \parencite{oficio-6507, oficio-5209}.

  \item[Actuaciones ante otros órganos] Comunicación dirigida a los tribunales
  \parencite{oficio-16156}; certificado de posesión efectiva ante el Registro Civil
  \parencite{oficio-14006}; boletas de honorarios cuando el deudor fallece durante el
  procedimiento \parencite{oficio-18305}.

  \item[Cajas de compensación] Aplicación de la \LIR{} a sus modificaciones
  \parencite{oficio-23665}.
\end{description}

\subsection{Oficios circulares}

Sólo uno se mantiene vigente: el Oficio Circular SIR \Nro 5, que refunde a los
anteriores y fija pautas para la correcta sustanciación de las audiencias del
procedimiento de renegociación \parencite{oficio-circular-5}. Su lectura es obligada
antes de comparecer a una audiencia ante la \superir{} (capítulo
\ref{cap:tram-renegociacion}).

\subsection{Resoluciones exentas de interés sistemático}

Además de las que aprueban instructivos y normas de carácter general, revisten interés
las que regularon el régimen transitorio de la reforma: la \resex{6.643}{2023} y la
\resex{6.644}{2023}, sobre cumplimiento de los artículos segundo y tercero transitorios
de la \Lreforma{} \parencite{resex-6643, resex-6644}, y las que actualizaron las nóminas
de veedores y liquidadores \parencite{resex-6638, resex-6640}. También la que aclaró la
vigencia de los oficios dictados durante la crisis sanitaria \parencite{resex-7101}.

%==============================================================================
\bloquePractico
%==============================================================================

\subsection{Cómo verificar la vigencia de una norma administrativa}

\begin{checklist}[Protocolo de verificación]
  \item Consultar la sección «Normas vigentes» del sitio de la \superir{}, que sólo
  publica las que conservan vigencia.
  \item Verificar si la norma fue modificada, complementada o derogada parcialmente por
  una resolución exenta posterior.
  \item Tratándose de instructivos, identificar la resolución aprobatoria y su fecha,
  dada la reiteración de la numeración.
  \item Revisar la sección «Normativa en trámite» para anticipar cambios inminentes
  (capítulo \ref{cap:normativa-tramite}).
  \item Contrastar la norma con los oficios de instrucción general posteriores, que
  pueden haber precisado su aplicación.
\end{checklist}

\begin{practica}[Fecha de corte]
Todo escrito que invoque normativa administrativa debería consignar la fecha en que se
verificó su vigencia. En un ámbito donde la producción normativa es mensual, la
diligencia profesional se acredita documentando cuándo se consultó.
\end{practica}

\subsection{Correspondencia entre procedimiento y normativa aplicable}

\begin{table}[htbp]
\centering
\small
\caption{Normativa administrativa aplicable a cada procedimiento}
\label{tab:normativa-por-procedimiento}
\begin{tabularx}{\textwidth}{@{}>{\raggedright\arraybackslash}p{0.26\textwidth} X@{}}
\toprule
\textbf{Procedimiento o actuación} & \textbf{Instrumentos aplicables} \\
\midrule
Reorganización ordinaria & NCG 4 (modelo de solicitud), NCG 19 (certificado y declaración
jurada), NCG 20 (informe de cumplimiento), Instructivo 3 (veedores) \\
\addlinespace
Reorganización simplificada & NCG 23 (declaración jurada), NCG 24 (modelo de propuesta) \\
\addlinespace
Renegociación de la Persona Deudora & NCG 28, Oficio Circular SIR 5 (audiencias) \\
\addlinespace
Liquidación simplificada & NCG 22 (declaración jurada), NCG 25 (plataformas
electrónicas), Res. Ex. 7.338 (citación a juntas), Oficio 9.378 \\
\addlinespace
Incautación y entrega de bienes & Instructivo 5, con la derogación parcial de la
Res. Ex. 16.492 \\
\addlinespace
Desvinculación laboral & Instructivo 6 (finiquitos electrónicos), Oficios 5.080, 334,
15.795, 1.765 y 6.499 \\
\addlinespace
Nominación de auxiliares & NCG 17 (sorteo), NCG 26 (requisitos e indicadores), NCG 16
(exámenes), NCG 10 (garantía de fiel desempeño) \\
\addlinespace
Juntas de acreedores & NCG 15 (juntas telemáticas) \\
\addlinespace
Cuentas y su objeción & NCG 7, NCG 27 (cuentas provisorias), NCG 18 (objeción de cuenta
final) \\
\addlinespace
Publicidad y notificaciones & NCG 14 \\
\addlinespace
Deberes de información del auxiliar & Instructivo 4 (hechos relevantes) \\
\bottomrule
\end{tabularx}
\end{table}
